\chapter[Kiểm thử và trạng thái hoàn thành]{Kiểm thử và trạng thái\\ hoàn thành}

\section{Phạm vi kiểm thử}

Hoạt động kiểm thử được chia thành hai nhóm. Kiểm thử kỹ thuật xác nhận mã nguồn có thể phân tích cú pháp và các frontend có thể tạo bản build. Kiểm thử nghiệp vụ tập trung vào các quy tắc cốt lõi của tài khoản, voucher, kiểm duyệt, mua hàng, phát hành E-Voucher và xác thực sử dụng.

\section{Kiểm thử kỹ thuật}

\begin{longtable}{|p{1.8cm}|p{3.5cm}|p{5.8cm}|p{2.3cm}|}
\hline
\textbf{Mã test} & \textbf{Chức năng} & \textbf{Dữ liệu/Điều kiện và kết quả mong đợi} & \textbf{Trạng thái} \\ \hline
\endfirsthead
\hline
\textbf{Mã test} & \textbf{Chức năng} & \textbf{Dữ liệu/Điều kiện và kết quả mong đợi} & \textbf{Trạng thái} \\ \hline
\endhead
TC-TECH-01 & Build frontend Khách hàng & Chạy \texttt{npm run build}; quá trình build hoàn tất và sinh thư mục đầu ra & Đạt, có cảnh báo \\ \hline
TC-TECH-02 & Build frontend Đối tác & Chạy \texttt{npm run build}; quá trình build hoàn tất và sinh thư mục đầu ra & Đạt, có cảnh báo \\ \hline
TC-TECH-03 & Build frontend Quản trị viên & Chạy \texttt{npm run build}; build hoàn tất; cảnh báo chunk lớn không làm dừng quá trình & Đạt, có cảnh báo \\ \hline
TC-TECH-04 & Kiểm tra cú pháp backend & Chạy \texttt{node --check} cho entry point và service chính; không phát sinh lỗi cú pháp & Đạt \\ \hline
TC-TECH-05 & Backend regression tests & Chạy \texttt{npm test}; 6 test bảo mật và nhất quán dữ liệu & Đạt \\ \hline
TC-TECH-06 & Frontend lint & Customer không còn lỗi; Admin còn 25 lỗi; Partner còn 40 lỗi & Chưa đạt \\ \hline
TC-TECH-07 & Dependency audit & Chạy \texttt{npm audit --omit=dev}; không còn vulnerability & Đạt \\ \hline
\caption[]{Kết quả kiểm thử kỹ thuật}\\
\end{longtable}

\section{Kiểm thử nghiệp vụ}

\footnotesize
\begin{longtable}{|p{1.4cm}|p{2.6cm}|p{3.6cm}|p{3.7cm}|p{1.5cm}|}
\hline
\textbf{Mã test} & \textbf{Chức năng} & \textbf{Dữ liệu/Điều kiện kiểm thử} & \textbf{Kết quả mong đợi} & \textbf{Trạng thái} \\ \hline
\endfirsthead
\hline
\textbf{Mã test} & \textbf{Chức năng} & \textbf{Dữ liệu/Điều kiện kiểm thử} & \textbf{Kết quả mong đợi} & \textbf{Trạng thái} \\ \hline
\endhead
TC-BIZ-01 & Đăng nhập hợp lệ & Tài khoản hoạt động, đúng tên đăng nhập và mật khẩu & Trả token và thông tin vai trò; cho phép truy cập đúng giao diện & Đạt \\ \hline
TC-BIZ-02 & Đăng nhập sai mật khẩu & Tài khoản tồn tại, mật khẩu không đúng & Từ chối đăng nhập và trả thông báo phù hợp & Đạt \\ \hline
TC-BIZ-03 & Tạo voucher hợp lệ & Đối tác hợp lệ; giá bán nhỏ hơn giá gốc; ngày và số lượng hợp lệ & Tạo voucher và lưu ở trạng thái phù hợp để gửi duyệt & Đạt \\ \hline
TC-BIZ-04 & Tạo voucher có giá bán không hợp lệ & Giá bán lớn hơn hoặc bằng giá gốc & Từ chối lưu và thông báo vi phạm quy tắc giá & Đạt \\ \hline
TC-BIZ-05 & Quản trị viên duyệt voucher & Voucher ở trạng thái Pending; Quản trị viên có quyền hợp lệ & Chuyển voucher sang Approved và cho phép công bố & Đạt \\ \hline
TC-BIZ-06 & Khách hàng xem voucher & Tồn tại voucher Approved và voucher chưa Approved & API/giao diện công khai chỉ trả voucher Approved còn hiệu lực & Đạt \\ \hline
TC-BIZ-07 & Đối tác redeem mã hợp lệ & Mã E-Voucher tồn tại, còn hạn, trạng thái Unused và đúng Đối tác/chi nhánh & Chuyển trạng thái sang Used và ghi nhận thời điểm sử dụng & Đạt \\ \hline
TC-BIZ-08 & Đối tác redeem mã đã dùng & Mã E-Voucher có trạng thái Used & Từ chối xác nhận sử dụng lần hai & Đạt \\ \hline
TC-BIZ-09 & Checkout và phát hành E-Voucher & Giỏ hàng hợp lệ, voucher Approved còn tồn kho, thanh toán qua callback đã xác minh & Đơn chuyển sang Paid và sinh đủ E-Voucher theo số lượng mua & Chưa xác nhận bằng kiểm thử tích hợp \\ \hline
TC-BIZ-10 & Admin xác nhận thanh toán thủ công & Đơn Pending và Quản trị viên hợp lệ & Cập nhật Paid thông qua service chung và phát hành E-Voucher tương ứng & Đạt \\ \hline
\caption[]{Kết quả kiểm thử nghiệp vụ cốt lõi}\\
\end{longtable}
\normalsize

Các kết quả nghiệp vụ trên được đối chiếu với phiên bản hiện thực hiện tại, log build, kiểm tra cú pháp backend và các ràng buộc trong database.

\section{Ma trận truy vết yêu cầu và kiểm thử}

Bảng dưới đây liên kết các nhóm yêu cầu chính với module hiện thực và nhóm test tương ứng. Mục tiêu của ma trận là chứng minh các phần được mô tả trong báo cáo có điểm đối chiếu rõ ràng trong mã nguồn và kế hoạch kiểm thử.

\begin{longtable}{|p{3.4cm}|p{4.4cm}|p{5.0cm}|}
\hline
\textbf{Nhóm yêu cầu} & \textbf{Module/Thành phần hiện thực} & \textbf{Nhóm kiểm thử đối chiếu} \\ \hline
\endfirsthead
\hline
\textbf{Nhóm yêu cầu} & \textbf{Module/Thành phần hiện thực} & \textbf{Nhóm kiểm thử đối chiếu} \\ \hline
\endhead
Tài khoản và phân quyền & \texttt{auth}, middleware JWT, hồ sơ Customer/Partner/Admin & Đăng ký, đăng nhập, quên mật khẩu, khóa tài khoản, kiểm tra quyền truy cập Admin/Partner \\ \hline
Đăng ký và duyệt đối tác & \texttt{auth}, \texttt{admin}, \texttt{Partners}, \texttt{Branches} & Đăng ký đối tác, trạng thái \texttt{Pending}, Admin duyệt/từ chối, Partner chỉ đăng nhập khi \texttt{Approved} \\ \hline
Quản lý voucher & \texttt{partner}, \texttt{admin/Voucher}, \texttt{Vouchers}, \texttt{Voucher\_Branches} & Tạo voucher hợp lệ/không hợp lệ, gửi duyệt, Admin duyệt/từ chối/tạm ẩn, Customer chỉ thấy voucher \texttt{Approved} \\ \hline
Mua hàng và thanh toán & \texttt{customer/order}, \texttt{Orders}, \texttt{Order\_Items} & Validate giỏ hàng, tạo đơn, thanh toán mô phỏng/sandbox, hủy/thất bại/quá hạn, Admin xác nhận thanh toán thủ công \\ \hline
Phát hành và sử dụng E-Voucher & \texttt{customer/order}, \texttt{partner}, \texttt{E\_Vouchers} & Chỉ phát hành sau \texttt{Paid}, mã duy nhất, QR/mã hiển thị cho khách, Partner redeem đúng đối tác/chi nhánh, không dùng lại mã đã \texttt{Used} \\ \hline
Đánh giá và khiếu nại & \texttt{Reviews}, \texttt{customer/complaint}, API xử lý khiếu nại Admin & Chỉ đánh giá sau mua, tạo khiếu nại, Admin phản hồi và cập nhật trạng thái \\ \hline
Báo cáo và vận hành & Dashboard Admin/Partner, \texttt{System\_Logs}, SSE event & Kiểm tra dashboard theo vai trò, dữ liệu không lộ giữa đối tác, cập nhật trạng thái và log vận hành \\ \hline
\caption[]{Ma trận truy vết yêu cầu, module và kiểm thử}\\
\end{longtable}

\section{Checklist yêu cầu}

Theo file \texttt{REQUIREMENTS\_CHECKLIST.md} cập nhật gần nhất, trạng thái tổng hợp là:

\begin{table}[H]
\centering
\begin{tabular}{|p{7cm}|p{3cm}|}
\hline
\textbf{Nhóm trạng thái} & \textbf{Số lượng} \\ \hline
Đã hoàn thành hoặc có thể đối chiếu & Hầu hết yêu cầu chức năng, dữ liệu và phi chức năng trong phạm vi demo \\ \hline
Nằm ngoài phạm vi package hiện tại & Video demo nếu giảng viên yêu cầu; mã nguồn ứng dụng được quản lý tại thư mục \texttt{application} và đóng gói riêng khi chốt bản nộp \\ \hline
\end{tabular}
\caption{Tổng hợp checklist yêu cầu}
\end{table}

\section{Các phần đã đáp ứng tốt}

\begin{itemize}[leftmargin=1.2cm]
    \item Hệ thống thể hiện rõ ba vai trò Khách hàng, Đối tác và Quản trị viên.
    \item Cơ sở dữ liệu quan hệ, ERD và dữ liệu mẫu hỗ trợ đối chiếu nghiệp vụ.
    \item Luồng Đối tác tạo voucher và Quản trị viên duyệt voucher đã được hiện thực.
    \item API công khai chỉ hiển thị voucher đã được phê duyệt.
    \item Luồng checkout mô phỏng, phát hành E-Voucher và Đối tác tra cứu/xác nhận sử dụng E-Voucher đã có phiên bản hoạt động.
    \item Dashboard và báo cáo cơ bản đã có cho Đối tác và Quản trị viên.
\end{itemize}

\section{Các điểm kiểm soát bàn giao}

\begin{itemize}[leftmargin=1.2cm]
    \item Script bảo trì dữ liệu trong thư mục database script dùng để kiểm tra/backfill E-Voucher khi database demo có đơn Paid cũ chưa đủ mã.
    \item Bằng chứng demo được chốt sau khi dữ liệu Supabase và cấu hình môi trường ổn định.
    \item Mã nguồn ứng dụng được đóng gói từ thư mục \texttt{application} khi bản cuối cùng ổn định.
\end{itemize}

\section{Rủi ro còn lại}

\begin{table}[H]
\centering
\begin{tabularx}{\textwidth}{|p{4.0cm}|Y|}
\hline
\textbf{Rủi ro} & \textbf{Hướng kiểm soát} \\ \hline
Dữ liệu demo cũ có thể chưa nhất quán với logic mới & Script kiểm tra/backfill E-Voucher khóa mã của đơn không còn Paid và bổ sung mã cho đơn Paid còn thiếu. \\ \hline
Thanh toán là mô phỏng/sandbox & Trình bày rõ đây là phạm vi học thuật; không dùng giao dịch tiền thật. \\ \hline
Frontend tách thành nhiều ứng dụng & Chuẩn hóa biến môi trường, API base URL và quy trình chạy demo cho từng vai trò. \\ \hline
Chưa có test tự động đầy đủ & Tiếp tục bổ sung unit/integration test cho checkout, callback thanh toán, redeem voucher và phân quyền API sau giai đoạn demo. \\ \hline
Dữ liệu môi trường cần được làm sạch trước khi nộp & Không đưa file môi trường thật, mật khẩu ứng dụng, API key hoặc chuỗi kết nối database thật vào gói nộp/public repository. \\ \hline
\end{tabularx}
\caption{Rủi ro và hướng kiểm soát}
\end{table}

\section{Giới hạn hiện tại}

Trong phạm vi đồ án, hệ thống tập trung vào việc chứng minh nghiệp vụ thương mại điện tử voucher và khả năng demo end-to-end. Một số giới hạn hiện tại gồm:

\begin{itemize}[leftmargin=1.2cm]
    \item Thanh toán sử dụng cơ chế mô phỏng/sandbox, chưa tích hợp quy trình đối soát và xử lý tiền thật.
    \item Test tự động chưa bao phủ đầy đủ toàn bộ luồng nghiệp vụ; phần lớn kiểm thử hiện tại dựa trên test plan, build và kiểm tra thủ công.
    \item Dữ liệu demo cần được chốt trước khi bảo vệ để tránh trường hợp đơn hàng cũ hoặc E-Voucher cũ không phản ánh đúng logic mới.
    \item Ba frontend được triển khai riêng theo vai trò nên cần cấu hình chính xác API base URL và biến môi trường khi demo.
\end{itemize}
